\chapter{Covers, isogenies and other transfers}\label{sec:transfer}

The routes in this chapter move the instance to another representation
(an isogenous curve, a curve at another level of the isogeny volcano, a
Jacobian, a Weil descent) in the hope that index calculus is cheaper there.
These routes produced the programme's clearest negatives, and one lemma
that explains most of them.

\begin{lemma}[transport conservation,
\art{docs/ic/RESEARCH_TRANSFER_AWARE_IC_20261010.md} \S1]\label{lem:conservation}
Let $\varphi:E\to E'$ be an $\FF_q$-rational isogeny (or any homomorphism
of rational points), $\kappa=\ker\varphi\cap E(\FF_q)$, $R\in E(\FF_q)$,
$\Fb'\subset\varphi(E(\FF_q))$ a factor base on the codomain and
$\Fb=\varphi^{-1}(\Fb')\cap E(\FF_q)$. Then the number of ordered
$m$-decompositions of $R$ over $\Fb$ is $\abs\kappa^{m-1}$ times the number
of ordered $m$-decompositions of $\varphi(R)$ over $\Fb'$.
\end{lemma}
\begin{proof}
Images of decompositions of $R$ are decompositions of $\varphi(R)$.
Conversely, for $\sum Q_i=\varphi(R)$ choose rational preimages $P_i$ of
the $Q_i$; then $\sum P_i-R\in\kappa$, and of the $\abs\kappa^m$ choices
exactly $\abs\kappa^{m-1}$ sum to $R$ (translate $P_1$ by the kernel
element).
\end{proof}

So a homomorphic transfer changes nothing but three curve-specific inputs:
the rational kernel symmetry $\abs\kappa^{m-1}$ (bounded by the rational
torsion; trivial on prime-order curves), the factor-base geometry
available on the destination, and the cost of arithmetic there. Sequences
of isogenies compose to one homomorphism; a loop is an endomorphism and
acts on the subgroup as a known scalar, useful only as a known-log
multiplier. Multi-curve factor bases glued by isogenies are single-curve
bases in disguise.

\paragraph{Conductor gaps in the ECC2K-130 class.} Jao, Miller and
Venkatesan~\cite{jmv2005} prove random self-reducibility among curves with
the same endomorphism ring, leaving open whether curves whose conductors
differ by a large prime could differ in hardness. The ECC2K-130 class has
conductor gap 263; we computed its structure
(\art{docs/ic/RESEARCH_ISOGENY_CONDUCTOR_GAP_20261007.md}) and ran
several rounds of experiments. Native counted-$S$ index calculus on two fixed CM
pairs gave floor/crater ratios $0.975$ (99\% $[0.603,1.451]$) at
$\ell=8209$ and $1.127$ $[0.922,1.342]$ at $\ell=131101$: no level effect
(192 verified records, 96 exact replays;
\art{research/ecbench_conductor_gap_20261012/}). Two-torsion
symmetrisation is a quotient isogeny: at a matched base of abscissae it
cut F4 row-reduction work $4.1$--$4.3\times$ ($\abs H=2$) and
$17.6$--$19.6\times$ ($\abs H=4$) over 1{,}128 systems. A correction
the next day showed that at a matched number of \emph{columns} the saving
disappears; it was the cost of a redundant orbit base. A combined Weil
descent used $234$--$529\times$ the F4 operations of the blind system
(\art{research/isogeny_conductor_gap_round3_20261011/}). The later rounds and
the full correction trail are in \S\ref{sec:iso-conductor}.

\paragraph{Fibre-aware index calculus.} Isogeny fibres are torsors whose
arithmetic is controlled by one cohomology class
(\art{RESEARCH_ISOGENY_FIBERS_ECDLP.md}); relation generation is
isogeny-equivariant at zero lifting cost, and a fibre-closed pair table was
the best lever of the conductor-gap rounds ($12.8\times$ rho in its best
cell). Measured against matched rho across sizes at fixed cofactor, the
envelope IC/rho ratio grows from $8.4$--$8.8$ to $35.1$--$37.0$ ($h=4$) and
from $13.4$--$20.1$ to $60.4$--$63.1$ ($h=72$), as $\#E^{1/6}$: the gap
widens, it does not close (816 audited records;
\art{research/fiber_ic_scaling_20261011/}).

\paragraph{Other structure.} Symmetrised $u$-frame bases saved
$0.84\times$ at $n=17$ and $0.76$--$0.77\times$ at $n=23$, but cost
$4.01\times$ at $n=31$ (\art{research/koblitz_symmetrised_e2e_20260927/}).
GHS and hyperelliptic covers: ECC2K-130 does map to hyperelliptic curves,
but the genus that would pay is missing. In odd characteristic the cubic
GHS transfer from $E_0(\FF_{q^3})$ to a genus-3 Jacobian now runs end to end
with double-large-prime index calculus, at about $q^{1.4}$ against rho's
$q^{3/2}$, within $1.15\times$ of a measured negation-map rho at
$q=2^{20}$ (\S\ref{sec:ladder}). The
isogeny-loop GLV/GLS work produced a measured $1.09$--$1.15\times$ faster
four-dimensional scalar multiplication (\art{research/isogeny_loop_glv_20261012/})
and a $\sqrt{\text{\'elu}}$ kernel cost exponent $0.822$ $[0.805,0.838]$;
these speed up curve arithmetic, not the discrete logarithm.

\paragraph{Covers, 21--24 October.} The Joux--Vitse genus-3 cover of
$E(\FF_{p^6})$ was, on 24 October, within $1.6$--$2.0\times$ of a pooled rho
constant at $p=1009$ (a lower bound, superseded in November by a matched
negation rho, \S\ref{sec:jac-jv}); measuring the per-trial cost on a genus-255 binary Jacobian
turned the $N=130$ GHS counterfactual's modelled margin negative ($-1.83$
bits); and two censuses bound what covers exist at toy sizes (genus floor 13
for odd-order $E/\FF_{q^4}$; a $2^{-l}$ genus-3 fraction over
$\FF_{2^{3l}}$). All are in \S\ref{sec:ladder3}. None touches a prime field
or a prime-degree binary field, where the GHS genus is out of reach.

\paragraph{The published scale of high-genus index calculus.} The largest
Weil-descent hyperelliptic discrete logarithms in the literature are at
genus 31. Jacobson, Menezes and Stein~\cite{jms2001} applied GHS to elliptic
curves over $\FF_{2^{62}}$, $\FF_{2^{93}}$, $\FF_{2^{124}}$ and
$\FF_{2^{155}}$ and obtained genus-31 curves $C_{62}$, $C_{93}$, $C_{124}$
and $C_{155}$ over $\FF_{2^2}$ to $\FF_{2^5}$, each with $\#J=2r$, $r$
prime. They solved $C_{62}$, $C_{93}$ and $C_{124}$ with Enge--Gaudry
(relation collection for $C_{124}$ took 379 CPU-days on a mixed cluster) and
only estimated $C_{155}$, at 26{,}290 days on one 1\,GHz Pentium III.
Velichka~\cite{velichka2008} solved $C_{155}$ ($\#J\approx2^{155}$) with
Vollmer's algorithm and sieving: 136{,}533 relations, about 1{,}735 CPU-days
in total, 21 days wall-clock, of which 14 days were linear algebra. No genus
above 31 appears in that thesis; the journal version~\cite{vjs2014} was not
available to us and is unchecked. Our genus-255 per-trial measurement
(\S\ref{sec:ladder3}) is a cost sample, not a solved instance, at about
eight times the largest published genus. The cover ECC2K-130 would need, of
genus about $2^{130}$ (\S\ref{sec:ecc2k130}), lies about $2^{125}$ times
beyond it.

\paragraph{Genus three against the strongest rho.} Rounds three and four
of the genus-3 work ran without a rho arm; a 23 October study supplied one
(\art{research/ghs3_vs_rho_20261023/}; it has a frozen protocol and tables
but no README, and no review is recorded in it). On round 4's instances,
$q=2^{14}$ to $2^{20}$, three curves per size, in round 4's unit and binary,
the reference was an ideal negation rho priced with Karatsuba products and
32 walks per inversion ($R_{\mathrm{ideal}}$, the primary reference, the
most favourable to rho) and a measured negation walk ($R_{\mathrm{meas}}$).
The frozen rule declares ``IC below the strong rho'' at a size only if every
curve is below both references at every weight of the band. At $q=2^{16}$
it is: IC/$R_{\mathrm{ideal}}$ is $0.70$, $0.71$, $0.70$ at $w=0$ and $0.64$
and $0.38$ (geometric means) at $w=0.1$ and $1$. At the other six sizes the
decision is \emph{unresolved}, with geometric means $0.97$--$1.52$ ($1.27$ at
$2^{20}$). Three things qualify the result. The registered unit prices one
subgroup multiplication as one $\FF_q$ multiplication ($c_\ell=1$); the
protocol's own sensitivity, $c_\ell=9$, puts every instance at $2.7$--$12.4\times$
rho. The ratio carries $\sqrt c$ of the cofactor, because IC depends on $q$
alone and rho on $\ell=\#E_0/c$: the three curves at $2^{16}$ all have
cofactor 4, and a post-hoc fit gives exponent $0.474$ in $c$ (theory $0.5$)
and $-0.081$ per bit of $q$ at fixed $c$ (theory $-0.167$). And the pooled
trend over all 21 instances is flat ($0.034$ $[-0.047,0.105]$ per bit), so
no crossover size is determined. This is the programme's only
preregistered decision of index calculus below the strongest rho. It is a
measurement of a known pipeline in a structured, weak class
($E_0/\FF_{q^3}$ with $4\mid\#E_0$, never prime order), where theory
predicts $q^{4/3}$ against $q^{3/2}$; it depends on the pricing of one
operation; and it is not yet reviewed. It updates the ledger paper's
``$1.15\times$ at $2^{20}$'' and bears on that paper's prediction of a
crossover at $q\le2^{22}$ once rho is measured: below rho at one size,
unresolved at the others, with no measurable pooled trend.

\paragraph{The same family through a plane quartic.} The $m=2$ type of the
cubic GHS cover gives a smooth plane quartic over $\FF_q$ directly. On a
plane quartic, Diem's index calculus is $\tilde O(q)$ in theory, against
$q^{4/3}$ for the hyperelliptic genus-3 route above. Before 18 October the
repository's plane-quartic ``pure IC'' toys could not be priced: their
logarithms came from enumerating the subgroup, and the degree-3 class keys
merged distinct classes. A canonical $\mathrm{Pic}^0$ arithmetic now gives
every class a unique key (checked against brute-force class counts on 21
curves and full Cayley tables at $q=5$). A generic addition costs 3{,}523
multiplications and 44 inversions, flat in $q$ up to $2^{31}-1$
(\art{research/quartic_pic_canonical_20261018/}).

On that arithmetic, a pipeline with relations on lines, a reduced factor base
with one large prime, Lanczos modulo $\ell$ and canonical descent solved and
verified 51 of 51 units. It ran on the GHS-weak family $E_0/\FF_{q^3}$ with
$4\mid\#E_0$, for $q=223$ to $262{,}111$ (\art{research/quartic_n4_20261019/};
Table~\ref{tab:prog_transfer-quartic}). Its reference is a matched
\emph{plain} rho, without the negation map, on the same subgroup, and the
order holds across the pricing band for every $q\ge4051$.
This was a preregistered decision of index calculus below a matched
\emph{plain} rho. Its successor, N4~v2
(\art{research/quartic_n4v2_20261020/}, 20--25 October), took the same
family and instances from $q=2^8$ to $2^{22}$ (63 instances, 501 index-calculus
units and 2{,}442 rho runs, all verified) and added three changes: double
large primes (D), structured Gaussian elimination with Lanczos (L), and a
negation-map rho. DL's exponent is the $\tilde O(q)$ of Diem's
analysis~\cite{diem2006,dt2008}. At $q=4{,}194{,}271$ one DL solve took
$12.3$ CPU-seconds against 616 for negation rho.

\begin{table}[t]
\centering\small
\begin{tabularx}{\linewidth}{@{}ll>{\raggedright}p{1.9cm}>{\raggedright}p{2.0cm}>{\raggedright}p{2.3cm}>{\raggedright}p{2.2cm}>{\raggedright\arraybackslash}X@{}}
\toprule
Study & Arm & Over B0 & IC/rho & IC exponent & Rho exponent & Crossover \\
\midrule
N4 & & & $0.122$ $[0.10,0.16]$ & $1.083$ $[1.040,1.129]$ &
  $1.485$ $[1.437,1.525]$ & wall time: only above $q\approx3\times10^4$ \\
\addlinespace
N4~v2 & D & $0.081$ & & & & \\
 & DL & $0.042$ $[0.042,0.042]$ & $0.005$ $[0.004,0.008]$ &
  $0.983$ $[0.975,0.993]$; $0.988$ on $k=17..22$ & $1.543$ $[1.494,1.583]$ &
  $q^*=555$ $[389,752]$ (interpolated); $q^*\approx1{,}889$ in CPU time \\
 & DL1 & $0.022$ & & & & \\
\bottomrule
\end{tabularx}
\caption[]{Plane-quartic index calculus on the GHS-weak family
$E_0/\FF_{q^3}$, $4\mid\#E_0$. N4 against a matched plain rho at
$q=262{,}111$, exponents fitted on $q\ge4051$; N4~v2 against the frozen
pipeline B0 and against negation rho at $q=2{,}097{,}091$ (the top size).
IC/rho in $\FF_q$-multiplication equivalents. Exponents: fitted cost
exponent in $q$, over the full range unless stated. Crossover: in counted
units unless stated (the N4 crossover is in wall time).
D: double large primes; DL: D with structured Gaussian elimination and
Lanczos; DL1 is post hoc. Sources: \art{research/quartic_n4_20261019/},
\art{research/quartic_n4v2_20261020/}.}
\label{tab:prog_transfer-quartic}
\end{table}
The qualifications are substantial. The host was L0 (load average 22--60);
the intervals resample only three units per size at $k\ge17$; the
conversion constants move D's high-range exponent from $1.009$ to $1.048$;
the frozen B0 records do not replay at any of 11 replayed sizes; the
negation rho fails the calibration rule at three sizes; and there is no
validation file. The program's closing note calls the result ``our
measurements of a known result''
(\art{research/plane_quartic_program_20261018/CLOSING.md}). It covers the
GHS-weak family only; nothing here is about a generic curve, a prime-order
curve or a standard curve.

\paragraph{Smith's isogeny and the walk to it.} The other genus-3 type of the
census is hyperelliptic. Smith's $(2,2,2)$-isogeny takes it to a plane quartic
when his Propositions~2 and~4 pass. The walk study looked for curves on which
Smith's construction succeeds by walking 2-isogenies on the elliptic side
(\art{research/smith_elliptic_walk_20261018/}; Table~\ref{tab:prog_transfer-smith}).
The preregistered expectation of about $0.186$ for the pass fraction is
falsified, upward. A labelled post hoc analysis attributes the frozen fit's
$o(p)$ label to censored walks at the small primes, and reads the
$p\equiv3\pmod4$ exponent as $\tilde\Theta(p)$. At $p\equiv1\pmod4$ most
walks stay inside small trapped sets.

\begin{table}[t]
\centering\small
\begin{tabular}{@{}llll@{}}
\toprule
Primes & Walks & Quantity & Value \\
\midrule
$p=7$ to $31$ & exhaustive & pass fraction & $0.53$--$0.61$ (random curves: $0.1857$) \\
 & & neighbour & $0.74$--$0.98$, at $1.5$--$3.6$ steps \\
\addlinespace
$p=2^{10}$ to $2^{20}$ & sampled & walk-cost exponent & $-0.14$ $[-0.80,0.95]$ \\
$p\equiv3\pmod4$ & successful & walk-cost exponent & $1.13$ $[1.00,1.26]$ \\
\bottomrule
\end{tabular}
\caption[]{The 2-isogeny walk to curves on which Smith's construction
succeeds. Pass fraction: share of descended curves passing Smith's full
test. Neighbour: share of failing curves with a passing curve in their
2-isogeny component, and the mean distance to it. Walk-cost exponent: in
$p$; the first is the frozen fit, which labels the walk cost $o(p)$; the
second is post hoc, over successful walks only. Source:
\art{research/smith_elliptic_walk_20261018/}.}
\label{tab:prog_transfer-smith}
\end{table}
A second study, also open and not yet merged, maps divisor classes through
Smith's isogeny and runs the whole route at toy sizes. The program's closing
note lists what stays open (\art{research/plane_quartic_program_20261018/CLOSING.md}).

\paragraph{Joux--Vitse: two arms of one cover.} The Joux--Vitse cover of
$E(\FF_{p^6})$ (full rational 2-torsion, $\#E=4\ell$, a weak class) has been
measured with two solvers, and their comparisons with rho do not agree.
The full account, with the limitations and the history, is in
\S\ref{sec:jac-jv}. The \emph{sieve arm} at $p=1511$ ($\ell=2^{61.4}$)
costs $0.00647$ $[0.00545,0.00799]$ of a matched negation rho in group
additions and $0.0156$ $[0.0131,0.0192]$ in textbook $\FF_p$ multiplications:
rho is $155\times$ and $64\times$ dearer, and the order holds across the
band. An independent validator passed it (\texttt{VAL-20261103-c31c76};
\art{research/jv_sieve_negrho_20261102/}). It has no \texttt{pcc} claim
record, so under the proof-carrying-claims rule it is a validation-gate
result whose statement is an unverified assertion. It is the programme's
one validated measurement of index calculus below a matched rho; the
validator classes it as accounting, not an improvement, and it holds on the
weak class only. The \emph{plain Nagao arm} was $1.59$--$1.98\times$ a pooled
rho constant at $p=1009$ on 24 October, a lower bound
(\art{research/jv_cover_p_scaling_20261024/}). Against a matched negation rho
it is above rho at $p=1301$ and 1451 and unresolved at 1601 in group
additions, and above at all three primes in $\FF_p$ multiplications
(\art{research/jv_cover_negrho_rescore_20261101/}). A factor-base line search
takes it below rho in additions but not in multiplications; the modelled
crossings move to $p^*\approx1185$ (additions) and $1799$ (multiplications),
without a validation file. The roughly thirtyfold gap between the two arms
is unexplained.

\paragraph{Cover censuses.} The minimal cubic GHS genus over an isogeny
class follows $\#E(K)\bmod4$: 13 for odd order, 5 for $2\bmod4$, 3 for
$0\bmod4$, and genus 3 is the only value below rho's $q^{3/2}$
(\art{research/ghs_cubic_cover_census_20261012/}); at $n=5$ the floors are
33, 33, 25 and 5 by 2-torsion type (\art{research/ghs_cover_census_n5_20261013/}).
No deployed prime-field or prime-degree binary curve has a usable subfield;
only the RFC~2409 Oakley groups 3 and 4 are reachable, at genus 15--16, where
exponent models disagree in sign and the measured genus-16 relation phase
alone is 5.2 bits slower than rho (\art{research/cover_route_prime_binary_20261008/},
\art{research/ghs_oakley_extended_20261023/}). All 15{,}791 cover maps in the
registry were assessed: 165 binary maps have genus-3 arithmetic and the
15{,}626 odd-characteristic genus-2 covers lack a Jacobian implementation
(\art{research/cover_transfer_assessment_20261013/}). An inventory of about
90 cover and decomposition studies found every route on prime, Koblitz and
prime-degree binary curves closed, negative or constant-factor
(\art{research/cover_decomposition_gap_search_20261018/}). Smith's
$(2,2,2)$-isogeny route was reproduced natively (full success $0.18638$
against a theoretical $0.18574$ over $10^6$ curves) and its sign bias shrinks
like $1/\sqrt p$ (\art{research/smith_prop24_native_20261014/},
\art{research/smith_sign_scaling_20261016/}).

\paragraph{Volcano levels and generators.} A 403{,}168-run census found the
level effect equivalent in 33 of 35 comparisons and different in none
(\art{research/volcano_ecdlp_census_20261017/}), and a large-sample rho study
found the floor/crater ratio $0.998$ $[0.986,1.010]$
(\art{research/conductor_gap_rho_large_sample_20261018/}). The one real level
effect is the automorphism group: $\sqrt3$ on the $j=0$ crater, confirmed at
$1.86$--$1.92\times$ (\art{research/drastic_prime_conductor_gap_round3_20261012/}).
At $\ell=2$, floors give $5.8\%$ fewer relations than craters, and 99\% of the
contrast is explained by the base's cofactor classes
(\art{research/isogeny-volcano-cost-v2_20261017/},
\art{research/generator_effect_ic_20261019/}).

\paragraph{Other transfer lines.} Index calculus on genus-2 to -4
Jacobians over $\FF_p$ (the ADH--Gaudry attack, measured against rho on the
same Jacobian) read below a distinguished-point rho in early rounds, but the
reference was later found slow and the ratios are lower bounds, unranked
under the band (\art{research/notes/index-calculus/RESEARCH_HYPERELLIPTIC_IC_RHO.md});
later rounds are in review. A weak-curve atlas of 922 exact models and 36
published cases ranks covering-model and pairing signals but leaves the
selected-subgroup advantage unmeasured
(\art{research/weak_curve_leaderboard_20261012/}); three complete isogeny
volcano graphs (24{,}823 vertices) were built and their graph-wide IC
surveys halted at checkpoints (\art{research/full_volcano_graphwide_ic_20261017/}).
Weil restriction of ECC2K-130 to extension fields costs at least $2^{131}$,
$2^{70.19}$ times rho (cluster ``transfer-boundaries'').

\paragraph{Prime-field covers, 11--12 November.} One study carried an
elliptic logarithm from $E/\FF_p$ to an isogenous curve and to the Jacobians
of genus-2, -3 and -4 covers $C\to E$, verifying every hop, and measured each
destination against rho on $E$
(\art{research/chained_prime_transfer_20261111/}). Genus-2 index calculus
with the full degree-1 base costs $8.0$--$40.8\times$ matched rho at 0~GAE
for $p=211$--$1619$, and $47.11\times$ (genus 2) and $141.97\times$ (genus 3)
at $p=64{,}007$; amortising the factor-base logs leaves an online cost that
ties a memory-matched precomputed rho table. The measurements are in
\S\ref{sec:jac-chain}. Two follow-ups explain why no cover of this kind
helps: the cover's relations push forward to the degenerate
$[2]Q=Q+Q$ on $E$, and the complementary factor never receives the
logarithm (\S\ref{sec:obs-pushforward}). Every figure in these studies is
an unverified assertion (\texttt{[pcc:unverified]} at the source); none
has a validation file.

\paragraph{Cover images as factor bases.} A cover $\varphi:C\to E$ can
also be used only to choose the base: $\Fb=x\circ\varphi(V)$ for a
structured set $V$ of $u$-coordinates on $C$, with relations and linear
algebra on $E$. Three rounds on 12 November asked whether the summation
system restricted to $\varphi(V)$ gives F4 smaller matrices than the same
system with random coefficients (its null) or than PKM's tower base of the
same size (\art{research/cover_image_fb_20261112/README.md},
\art{research/cover_image_fb_r2_20261112/README.md},
\art{research/cover_image_fb_r3_m3_20261112/README.md}). The setting is
$p=2013265921$, Kummer towers, 3 planted and 3 random targets per cell,
cover maps of degree 2 to 4 (genus 2 to 5), with no timeouts and 0 planted
violations. Round~1's preregistered metric, the solving degree $D$, is
floored by the input degree ($S_3(u^2,\dots)$ has degree 8 in the bottom
tower variable), so its positive control failed on its letter and the
$D$-based verdict was not read; matrix width $W$ (the widest F4 matrix) was
read post hoc and preregistered for rounds 2 and~3. Table~\ref{tab:cover-image-fb}
gives the widths. At $m=2$ every map sits at its null in $W$: the
\texttt{quad}/null ratio is $0.97$--$1.02$ for $t=5$--$8$ ($1.31$ at $t=4$,
$0.81$ at $t=3$), \texttt{cubic} $0.95$--$1.01$, \texttt{generic3} $1.00$, and
\texttt{quart} $1.31$, $1.00$, $0.96$ at $t=5$, 6, 7, the $1.31$ being the cover
doing worse and the only cell outside the $1.25$ band; round~2 scored
\texttt{cubic} and \texttt{generic3} as the falsifier H0 and \texttt{quart} as
inconclusive. Every map is
$1.22$--$5.47\times$ wider than the tower at equal $\abs\Fb$, because substituting a degree-$d$ map raises $S_3$ to degree $4d$
in two summands. One $m=3$, $t=3$ cell had \texttt{cubic} and \texttt{quart}
narrower than the tower ($1{,}866$ and $1{,}876$ against $3{,}285$); round~3
tested it at $t=4$ with nulls, and all four maps landed on the falsifier H0
at $1.55$--$1.99\times$ the tower's width. At $t=3$ the cubic was $0.59\times$
its own null, at $t=4$ $1.00\times$. The $t=5$ cell exceeded the 16\,GB
memory limit. The study closes the line at these sizes ($m\le3$, $t\le8$,
small $p$); it reports no $S$, no relation-phase cost and no rho
comparison, and all values are \texttt{[pcc:unverified]}. PKM's tower was
itself already equal to its null (\S\ref{sec:pa-pkm}).

\begin{table}[t]
\centering\footnotesize
\begin{tabular}{@{}rrrrrrrrr@{}}
\toprule
$t$ & $\abs\Fb$ & PKM tower & \texttt{quad} & \texttt{quad} null & \texttt{generic} & \texttt{cubic} & \texttt{generic3} & \texttt{quart} \\
\midrule
4 & 16 & 280 & 645 & 491 & 565 & -- & -- & -- \\
5 & 32 & 606 & 1{,}629 & 1{,}601 & 1{,}594 & 1{,}663 & 1{,}749 & 3{,}315 \\
6 & 64 & 2{,}451 & 3{,}141 & 3{,}233 & 3{,}247 & 3{,}755 & 3{,}728 & 8{,}436 \\
7 & 128 & 5{,}031 & 6{,}159 & 6{,}355 & 6{,}330 & 17{,}126 & 17{,}313 & 16{,}453 \\
8 & 256 & 9{,}119 & 11{,}723 & 11{,}570 & 11{,}580 & -- & -- & -- \\
\midrule
\multicolumn{9}{@{}l}{$m=3$, $t=4$, $\abs\Fb=16$: width, ratio to tower, ratio to own null, $D$ (input degree)}\\
\multicolumn{9}{@{}l}{\quad PKM tower $12{,}314$, 1, $0.99$, 6 (4);\quad \texttt{quad} $24{,}482$, $1.99$, $1.00$, 8 (8);\quad \texttt{cubic} $19{,}911$, $1.62$, $1.00$, 12 (12);}\\
\multicolumn{9}{@{}l}{\quad \texttt{quart} $19{,}079$, $1.55$, $0.96$, 16 (16);\quad \texttt{generic3} $19{,}863$, $1.61$, $1.00$, 13 (13)}\\
\bottomrule
\end{tabular}
\caption[]{Cover-image factor bases: widest F4 matrix $W$ at $m=2$ (Kummer
tower control), and the $m=3$, $t=4$ cell. \texttt{quad}, \texttt{generic}:
degree 2, genus 2 and 3 (round~1); \texttt{cubic}, \texttt{generic3},
\texttt{quart}: degree 3, 3 and 4, genus 4, 5 and 5 (round~2). Rows $t=4$ and
$t=8$ of round~2 were not run. Sources:
\art{research/cover_image_fb_20261112/README.md},
\art{research/cover_image_fb_r2_20261112/README.md},
\art{research/cover_image_fb_r3_m3_20261112/README.md};
\texttt{[pcc:unverified]}, stage diagnostic.}
\label{tab:cover-image-fb}
\end{table}

\paragraph{Named obstructions.} Blocked requirements are kept open as
problem records (\art{problems/README.md}), among them the complete-cost
calibration of root-finding IC
(\art{problems/rootfinding-complete-cost-calibration.md}), with exit
criteria rather than abandoned.
